% TEMPLATE-MILC-v3.tex - plantilla maestra para todas las guias MILC v3
% Ver scripts/generadores/build-guia-milc.py para la lista completa de
% claves esperadas. El script valida que no queden placeholders sin
% reemplazar antes de escribir el archivo final.

\documentclass[11pt,letterpaper]{article}
\usepackage[margin=1.75cm]{geometry}
\usepackage[table]{xcolor}
\usepackage{fontspec}
\usepackage[spanish]{babel}
\usepackage{tikz}
\usepackage[most]{tcolorbox}
\usepackage{tabularx}
\usepackage{array}
\usepackage{fancyhdr}
\usepackage{titlesec}
\usepackage{ragged2e}
\usepackage{enumitem}
\usepackage{microtype}

\setmainfont{Helvetica Neue}
\setsansfont{Helvetica Neue}
\setlength{\parindent}{0pt}
\setlength{\parskip}{6pt}
\hyphenpenalty=200
\exhyphenpenalty=200
\pretolerance=400
\tolerance=2000
\setlength{\emergencystretch}{1em}
\renewcommand{\arraystretch}{1.25}
\setlist{leftmargin=5mm,itemsep=1mm,topsep=1mm}
\newcolumntype{Y}{>{\RaggedRight\arraybackslash}X}

\definecolor{milcMagenta}{HTML}{E6007E}
\definecolor{milcVino}{HTML}{5A0038}
\definecolor{milcTurquesa}{HTML}{0093A5}
\definecolor{milcVerde}{HTML}{58A923}
\definecolor{milcMostaza}{HTML}{E5B400}
\definecolor{milcOcre}{HTML}{B86600}
\definecolor{milcNegro}{HTML}{241816}
\definecolor{milcGris}{HTML}{F7F7F4}
\definecolor{milcLinea}{HTML}{E8E2DD}
\definecolor{milcGradePrimary}{HTML}{0066FF}
\definecolor{milcGradeDark}{HTML}{003D99}
\definecolor{milcGradeSoft}{HTML}{D6E8FF}
\definecolor{milcGradeAccent}{HTML}{CFE7FF}
\definecolor{milcGradeText}{HTML}{FFFFFF}
\color{milcNegro}

\pagestyle{fancy}
\fancyhf{}
\lhead{\small Tecnología e Informática · Grado 11}
\rhead{\small Guía 7 · MILC v3}
\cfoot{\small\thepage}
\renewcommand{\headrulewidth}{0pt}

\newtcolorbox{titlebox}[1]{
  enhanced,colback=#1,colframe=#1,arc=7mm,boxrule=0pt,
  left=7mm,right=7mm,top=3mm,bottom=3mm,
  before skip=8pt,after skip=8pt,
  fontupper=\sffamily\bfseries\Large\color{white}
}

\newtcolorbox{softbox}[3]{
  enhanced,colback=#2,colframe=#1,borderline west={4pt}{0pt}{#1},
  arc=2mm,boxrule=.4pt,left=4mm,right=4mm,top=3mm,bottom=3mm,
  before skip=6pt,after skip=8pt,title={#3},coltitle=white,
  colbacktitle=#1,fonttitle=\sffamily\bfseries,fontupper=\RaggedRight,
  attach boxed title to top left={xshift=3mm,yshift=-2mm},
  boxed title style={arc=2mm, boxrule=0pt}
}

\newtcolorbox{drawbox}[2]{
  enhanced,colback=white,colframe=#1,arc=4mm,boxrule=1pt,
  left=5mm,right=5mm,top=4mm,bottom=4mm,before skip=8pt,after skip=8pt,
  title={#2},coltitle=white,colbacktitle=#1,fonttitle=\sffamily\bfseries,
  fontupper=\RaggedRight,
  attach boxed title to top left={xshift=4mm,yshift=-2mm},
  boxed title style={arc=3mm, boxrule=0pt}
}

\newtcolorbox{infoband}[1]{
  enhanced,colback=#1!8,colframe=#1!50,arc=2mm,boxrule=.5pt,
  left=4mm,right=4mm,top=2mm,bottom=2mm,before skip=4pt,after skip=4pt,
  fontupper=\small\RaggedRight
}

% Puente narrativo entre secciones (contrato editorial Conéctate).
% Da continuidad: "Acabas de X. Ahora vas a Y porque Z."
% Visualmente: banda izquierda en color del grado, texto en itálica pequeña.
\newtcolorbox{puentebox}{
  enhanced,colback=milcGradeSoft,colframe=milcGradeDark,
  borderline west={3pt}{0pt}{milcGradeDark},
  arc=0mm,boxrule=0pt,left=5mm,right=4mm,top=2.5mm,bottom=2.5mm,
  before skip=4pt,after skip=8pt,
  fontupper=\itshape\small\color{milcNegro}\RaggedRight
}
\newcommand{\puente}[1]{%
  \begin{puentebox}%
    \textcolor{milcGradeDark}{\textbf{\textrightarrow}}\hspace{2mm}#1%
  \end{puentebox}%
}

\newcommand{\linead}[1]{\noindent\color{milcLinea}\rule{#1}{0.5pt}\color{milcNegro}}
\newcommand{\respuesta}[1]{\par\vspace{2mm}\foreach \n in {1,...,#1}{\noindent\color{milcLinea}\rule{\linewidth}{0.5pt}\par\vspace{5mm}}\color{milcNegro}}
\newcommand{\checkbox}{\textcolor{milcGradeDark}{\fbox{\phantom{x}}}\hspace{5pt}}

\newcommand{\phasecard}[4]{
\begin{tcolorbox}[enhanced,colback=#3,colframe=#2,arc=3mm,boxrule=.5pt,height=3.05cm,valign=center,halign=center,left=1.5mm,right=1.5mm,top=2mm,bottom=2mm]
{\sffamily\bfseries\fontsize{22}{24}\selectfont\textcolor{#2}{#1}}\par
{\sffamily\bfseries\small #4}
\end{tcolorbox}
}

\begin{document}
\thispagestyle{empty}
\begin{tikzpicture}[remember picture,overlay]
  \fill[white] (current page.south west) rectangle (current page.north east);
  \path[fill=milcGradePrimary,rounded corners=16mm]
    ([xshift=.55cm,yshift=-.55cm]current page.north west) rectangle
    ([xshift=-.55cm,yshift=3.65cm]current page.south east);

  \node[anchor=north west,text=milcGradeText] at ([xshift=2.0cm,yshift=-1.85cm]current page.north west)
    {\sffamily\bfseries\fontsize{14}{16}\selectfont GRADO};
  \node[anchor=north west,text=milcGradeText,opacity=.82] at ([xshift=2.0cm,yshift=-2.75cm]current page.north west)
    {\sffamily\bfseries\fontsize{9}{11}\selectfont SERIE GUÍAS · TECNOLOGÍA E INFORMÁTICA};

  \begin{scope}[shift={([xshift=-3.05cm,yshift=-2.55cm]current page.north east)}]
    \fill[white,opacity=.80] (-.62,-.52) rectangle (.62,.52);
    \draw[milcGradeText,line width=1pt] (-.62,-.52) rectangle (.62,.52);
    \fill[milcGradeDark] (-.42,-.38) rectangle (-.22,.06);
    \fill[milcGradeAccent] (-.10,-.38) rectangle (.10,.24);
    \fill[milcGradePrimary!60!black] (.22,-.38) rectangle (.42,.38);
  \end{scope}

  \node[anchor=west,text=milcGradeText] at ([xshift=1.9cm,yshift=-12.85cm]current page.north west)
    {\sffamily\bfseries\fontsize{124}{124}\selectfont 11};
  \draw[milcGradeText,line width=1.7pt]
    ([xshift=4.52cm,yshift=-11.68cm]current page.north west) circle (.13cm);

  \node[anchor=west,text=milcGradeText,text width=8.7cm,align=left] at ([xshift=10.35cm,yshift=-11.6cm]current page.north west)
    {
     {\sffamily\bfseries\fontsize{19}{22}\selectfont Imagen propia ---\\cómo se ve mi marca\\en foto y video}\\[4mm]
     {\sffamily\bfseries\fontsize{23}{26}\selectfont Guía 7-11-TIC}\\[2mm]
     {\sffamily\fontsize{13}{16}\selectfont Período 1 · Presencia y marca digital}\\[5mm]
     {\sffamily\bfseries\fontsize{11}{13}\selectfont Institución Educativa Sor María Juliana}\\[1mm]
     {\sffamily\fontsize{10}{12}\selectfont Autor: Dr. Álvaro Cárdenas Orozco}
    };

  \path[fill=milcGradeSoft,rounded corners=7mm]
    ([xshift=1.35cm,yshift=.85cm]current page.south west) rectangle
    ([xshift=-1.35cm,yshift=3.15cm]current page.south east);
  \draw[milcGradeDark,line width=.65pt,rounded corners=7mm]
    ([xshift=1.35cm,yshift=.85cm]current page.south west) rectangle
    ([xshift=-1.35cm,yshift=3.15cm]current page.south east);
  \node[anchor=center,text=milcNegro,text width=17.0cm,align=center] at ([yshift=2.0cm]current page.south)
    {\sffamily\fontsize{10}{12}\selectfont Metodología MILC · Apertura ancestral · Pensamiento computacional · Triángulo Dussel-Estoicismo-Floridi\\
    \textcolor{milcGradeDark}{\bfseries Producto:} 3 fotos de perfil profesional + 1 video corto ($\leq$60 s) de presentación};
\end{tikzpicture}
\mbox{}
\newpage

\section*{Apertura: Imagen propia --- cómo se ve mi marca en foto y video}

\begin{softbox}{milcGradeDark}{milcGradeSoft}{Saber ancestral}
El \textbf{retrato del oficio} era un género visual con propósito: el panadero con su pan, el pescador con su atarraya, la tejedora con su telar, el campesino con su azadón. No era \emph{selfie}. Era una imagen que decía sin palabras quién hacía qué, dónde, y con qué orgullo. La foto comunicaba dignidad, oficio y competencia antes de que la persona dijera su nombre.
\end{softbox}

\begin{softbox}{milcTurquesa}{milcGris}{Saber contemporáneo}
Hoy tu marca necesita lo mismo en formato digital: una \textbf{foto de perfil} que comunique tu oficio, una \textbf{foto de ambiente} que muestre dónde trabajas, y un \textbf{video corto} (60 segundos máximo) donde te presentas. En redes y perfiles, esa pieza visual habla por ti antes que tu texto. Quien te ve decide en menos de 3 segundos si quiere leerte.
\end{softbox}

\begin{softbox}{milcMostaza}{milcGris}{Pregunta-puente}
\textit{¿Qué hacía que una foto antigua de oficio comunicara dignidad y competencia sin decir una palabra? ¿Qué de eso puedes replicar hoy con un celular?}
\end{softbox}

\begin{softbox}{milcVerde}{milcGris}{Saber y saber hacer}
Al terminar podrás: (1) \textbf{identificar} qué señales visuales transmite una imagen profesional vs una casual; (2) \textbf{analizar} 3 fotos profesionales que admiras; (3) \textbf{crear} tus 3 fotos + 1 video con criterio; (4) \textbf{evaluar} la imagen propia de un compañero.
\end{softbox}

\small
\begin{tabularx}{\textwidth}{>{\bfseries}p{3.0cm}Y}
\rowcolor{milcGradePrimary!12}Autor & Dr. Álvaro Cárdenas Orozco\\
DBA & Comunicación profesional en entornos digitales (MEN, Lineamientos T\&I)\\
Referentes & Floridi (infoética) · Dussel (estética de la liberación) · Estoicismo\\
Producto final & 3 fotos de perfil profesional + 1 video corto ($\leq$60 s) de presentación\\
\end{tabularx}
\normalsize

\puente{Antes de empezar, mira el viaje completo. Hoy le pones \emph{cara} a tu marca --- no la selfie de tu cumpleaños, una imagen que comunique quién eres profesionalmente.}

\begin{titlebox}{milcGradeDark}
Ruta MILC: así vamos a aprender
\end{titlebox}
\begin{tabularx}{\textwidth}{>{\centering\arraybackslash}X>{\centering\arraybackslash}X>{\centering\arraybackslash}X>{\centering\arraybackslash}X}
\phasecard{1}{milcVerde}{milcGris}{Escuta\\[-1mm]\small Observo, escucho y pregunto.} &
\phasecard{2}{milcTurquesa}{milcGris}{Sistematización\\[-1mm]\small Pensamiento computacional.} &
\phasecard{3}{milcMagenta}{milcGris}{Praxis\\[-1mm]\small Construyo y aplico.} &
\phasecard{4}{milcVino}{milcGris}{Evaluación\\[-1mm]\small Reflexión liberadora.}
\end{tabularx}


\puente{Antes de tomarte tu foto, vas a leer otras. Las imágenes que funcionan tienen decisiones técnicas atrás: encuadre, fondo, mirada, iluminación. Entrenar el ojo es el primer paso.}

\begin{titlebox}{milcVerde}
Fase 1 · Escuta
\end{titlebox}

\begin{softbox}{milcVerde}{milcGris}{Escena para escuchar}
\textbf{Actividad 1 · IDENTIFICA --- Foto profesional vs foto casual} (15 min · individual). Busca en LinkedIn (o portafolios en línea) \textbf{3 fotos profesionales} de personas que admiras y \textbf{3 fotos casuales} (Instagram, TikTok) de cualquier persona. Para cada par: ¿qué diferencias notas en encuadre, fondo, mirada, luz?
\end{softbox}

\checkbox Encontré 3 fotos profesionales que comunican oficio o competencia.\\
\checkbox Las contrasté con 3 fotos casuales de mi propio feed.\\
\checkbox Identifiqué 4 diferencias técnicas concretas (encuadre/fondo/luz/intención).

\begin{infoband}{milcVerde}
\textbf{Lo que va al cuaderno (Actividad 1 · IDENTIFICA).} Título: «Actividad 1 --- Foto profesional vs casual». \textbf{Formato:} tabla de 5 columnas (Foto | Encuadre | Fondo | Luz | Qué transmite), 6 filas (3 profesionales + 3 casuales). \textbf{Sabes que terminaste cuando} cada celda nombra una decisión técnica concreta (no ``se ve mejor''), y al final puedes resumir en 2 renglones qué diferencia visual hace la categoría ``profesional''.
\end{infoband}

\puente{Ya viste qué hace fuerte una imagen de oficio moderna. Ahora vamos a aprender las \textbf{cuatro decisiones técnicas} que la diferencian de una foto casual: encuadre, fondo, luz, intención.}

\begin{titlebox}{milcTurquesa}
Fase 2 · Sistematización con pensamiento computacional
\end{titlebox}

\begin{softbox}{milcTurquesa}{milcGris}{Cuatro pilares aplicados al tema}
Toda imagen profesional sigue 4 decisiones técnicas. Vamos a descomponerlas con los pilares del pensamiento computacional para que puedas hacer las tuyas con criterio y un celular.
\end{softbox}

\small
\begin{tabularx}{\textwidth}{>{\bfseries\RaggedRight\arraybackslash}p{3.6cm}Y}
\rowcolor{milcTurquesa!90}\color{white}Pilar & \color{white}\bfseries Aplicación al tema\\
1. Descomposición & Una imagen profesional se descompone en 4 decisiones: (1) \textbf{encuadre} (plano medio para retrato, plano general para ambiente), (2) \textbf{fondo} (limpio, coherente, no distractor), (3) \textbf{luz} (frontal natural mejor que flash, evita contraluces fuertes), (4) \textbf{intención} (sabes qué transmitir antes de disparar).\\
2. Reconocimiento de patrones & Toda imagen sigue el patrón \textbf{regla de los tercios}: divide mentalmente la imagen en 9 partes y pon el sujeto en una de las líneas. El centro exacto se siente plano; la línea, equilibrado.\\
3. Abstracción & Lo esencial de tu imagen cabe en una pregunta: \emph{¿qué dirá esta foto si la ve alguien que no me conoce, en 3 segundos?}. Si la respuesta es ``algo cualquiera'', todavía no abstrajiste tu propuesta visual.\\
4. Algoritmo conceptual & Pasos: (1) decide qué quieres comunicar (oficio, ambiente, acción); (2) elige fondo y luz; (3) usa el temporizador del celular (no selfie con brazo); (4) toma 20+ fotos; (5) descarta el 95\%, queda con la mejor.\\
\end{tabularx}
\normalsize

\begin{drawbox}{milcTurquesa}{Anatomía de una foto profesional con celular}
\textbf{Retrato (perfil):} plano medio, mirada al lente o a 45°, fondo limpio, luz natural lateral. \\[1mm] \textbf{Ambiente:} plano más abierto que muestra dónde trabajas, sujeto a un tercio del encuadre. \\[1mm] \textbf{Acción:} foto tomada \emph{mientras} haces tu oficio (no posada), con un objeto característico. \\[1mm] \textbf{Edición:} ajustes mínimos (exposición + contraste). Sin filtros agresivos. Sin filtros de cara.
\end{drawbox}

\begin{softbox}{milcMostaza}{milcGris}{Errores comunes y su corrección}
(1) Selfie con brazo extendido: la distorsión hace ver mal hasta a un modelo. (2) Fondos caóticos (cuarto desordenado, mucha gente): pierdes el foco. (3) Filtros agresivos: rompen la coherencia con cualquier marca seria. (4) Una sola foto: nunca sabrás si pudiste hacerlo mejor.
\end{softbox}

\begin{infoband}{milcTurquesa}
\textbf{Lo que va al cuaderno (Actividad 2 · ANALIZA).} Título: «Actividad 2 --- Análisis de 3 imágenes que admiro». \textbf{Formato:} 3 fichas, una por foto, cada una con: (a) descripción de la foto en 1 frase, (b) las 4 decisiones técnicas (encuadre/fondo/luz/intención), (c) qué transmite. \textbf{Sabes que terminaste cuando} las 3 fichas tienen las 4 decisiones identificadas con precisión, no con generalidades.
\end{infoband}

\puente{Tienes el criterio. Toca producir: 3 fotos y 1 video corto con celular y un par de minutos. Sin presupuesto, sin sesión profesional --- con decisiones acertadas.}

\begin{titlebox}{milcMagenta}
Fase 3 · Praxis con pensamiento computacional
\end{titlebox}

\begin{softbox}{milcMagenta}{milcGris}{Construyendo el producto con los cuatro pilares}
\textbf{Actividad 3 · CREA --- Tus 3 fotos + 1 video corto} (40 min · individual o con un compañero que tome las fotos). Hora de producir. Vas a generar 3 fotos (retrato + ambiente + acción) y 1 video de presentación de máximo 60 segundos --- con celular, sin presupuesto, con criterio.
\end{softbox}

\small
\begin{tabularx}{\textwidth}{>{\bfseries\RaggedRight\arraybackslash}p{3.6cm}Y}
\rowcolor{milcMagenta!90}\color{white}Pilar & \color{white}\bfseries Aplicación al producto\\
Descomposición & Tu sesión se descompone en 4 piezas: retrato (perfil), ambiente, acción, video corto. Cada una con propósito distinto y decisión técnica distinta.\\
Patrones & Sigue el patrón ``20 fotos por foto final'': nunca te quedes con la primera. La diferencia entre amateur y profesional no es la cámara: es la cantidad de tomas que descartas.\\
Abstracción & Cada foto comunica una sola idea sobre ti. Si tu retrato muestra ``soy serio'' y tu video dice ``soy divertido'', la marca pierde coherencia. Decide \emph{antes} qué transmitir.\\
Algoritmo de armado & Algoritmo: (1) elige lugar y hora (luz natural, mañana o atardecer); (2) pide a un compañero o usa temporizador; (3) toma 20+ fotos por cada una de las 3 categorías; (4) graba 5 takes del video, queda con 1; (5) edición mínima en el celular; (6) sube al sitio o portafolio.\\
\end{tabularx}
\normalsize

\begin{drawbox}{milcMagenta}{Checklist antes de cerrar la sesión de fotos}
\checkbox 3 fotos finales seleccionadas (retrato + ambiente + acción). \\[1mm] \checkbox 1 video de presentación de 30-60 segundos, audio limpio. \\[1mm] \checkbox Cero fotos con brazo extendido (selfie clásica). \\[1mm] \checkbox Fondos limpios o intencionalmente narrativos. \\[1mm] \checkbox Edición mínima, sin filtros de cara ni colores agresivos. \\[1mm] \checkbox Coherente con tu sistema visual (paleta, tono) de la Guía 2.
\end{drawbox}

\begin{softbox}{milcMagenta}{milcGris}{Plantilla de trabajo}
\textbf{Retrato.} \textit{Descripción breve + decisiones técnicas usadas + URL/captura.} \\[1mm] \textbf{Ambiente.} \textit{Lo mismo.} \\[1mm] \textbf{Acción.} \textit{Lo mismo.} \\[1mm] \textbf{Video (60 s).} \textit{Guion de 3 ideas + URL del video subido.}
\end{softbox}

\begin{infoband}{milcMagenta}
\textbf{Lo que va al cuaderno (Actividad 3 · CREA).} Título: «Actividad 3 --- Mi imagen de marca · v1». \textbf{Formato:} ficha por cada uno de los 4 entregables (3 fotos + 1 video) con descripción y URL/captura. \textbf{Extensión:} una página de cuaderno. \textbf{Sabes que terminaste cuando} un extraño que mira los 4 entregables puede decir, en 1 frase, qué haces y cómo te ves haciéndolo.
\end{infoband}

\puente{Lo que sigue resume \emph{qué} entregas hoy y \emph{cómo se sabe} si tu imagen funciona. El criterio principal: que un extraño pueda mirarla 3 segundos y entender qué haces.}

\begin{titlebox}{milcMostaza}
Producto de la sesión
\end{titlebox}
\textbf{3 fotos profesionales + 1 video corto coherentes con tu marca}.
\begin{softbox}{milcMostaza}{milcGris}{Criterios de calidad}
(1) Las 3 fotos tienen propósito distinto (retrato / ambiente / acción). (2) Cero selfies con brazo extendido o filtros de cara. (3) El video tiene audio limpio y mensaje claro en máximo 60 segundos. (4) Las 4 piezas son coherentes con tu sistema visual (paleta, tono) de la Guía 2. (5) Un extraño puede entender qué haces en 3 segundos de mirada.
\end{softbox}

\puente{Antes de cerrar, mira tu imagen desde las cinco dimensiones humanas. La foto profesional toca temas profundos: cómo te ves, cómo te ven, qué expectativa construyes.}

\begin{titlebox}{milcVino}
Fase 4 · Evaluación liberadora — cinco dimensiones
\end{titlebox}

\begin{softbox}{milcVino}{milcGris}{Sentido de la evaluación}
Evaluar no es solo revisar una nota. Es reconocer qué aprendí, qué cambió en mí, qué emociones manejé, cómo me relaciono con la ciudad y la comunidad, y qué le dejaré a quien viene detrás.
\end{softbox}

\textbf{Mi reflexión liberadora en cinco dimensiones.} Completa cada frase iniciada:

\small
\begin{tabularx}{\textwidth}{>{\bfseries\RaggedRight\arraybackslash}p{3.4cm}Y>{\RaggedRight\arraybackslash}p{4.6cm}}
\rowcolor{milcVino}\color{white}Dimensión & \color{white}\bfseries Mi respuesta (frame starter) & \color{white}\bfseries Algo que descubrí\\
1. Personal & Lo que más cambió en mí fue \linead{6.5cm} & \linead{4.2cm}\\
2. Emocional & La emoción más fuerte fue \linead{2.5cm} y la regulé \linead{3cm} & \linead{4.2cm}\\
3. Ciudadana & En democracia esto importa porque \linead{6cm} & \linead{4.2cm}\\
4. Local (Cartago) & En mi barrio sería útil para \linead{6.5cm} & \linead{4.2cm}\\
5. Intergeneracional & A mi abuela/abuelo le diría \linead{6.5cm} & \linead{4.2cm}\\
\end{tabularx}
\normalsize

\begin{infoband}{milcVino}
\textbf{Para la próxima sustentación.} Vuelve a esta tabla en seis meses. Verás que las respuestas cambian — no porque lo aprendido cambie, sino porque tú cambias. La evaluación liberadora es una conversación contigo a través del tiempo.
\end{infoband}


\puente{Y para terminar, tres voces sobre la imagen pública. Dussel sobre la voz visual silenciada, Epicteto sobre lo que depende de ti, Floridi sobre la imagen como información en la infoesfera.}

\begin{titlebox}{milcGradeDark}
Triángulo de pensamiento · cierre filosófico
\end{titlebox}

\begin{softbox}{milcVino}{milcGris}{Dussel · lente del nosotros}
\textit{``La belleza estética nace donde una voz silenciada encuentra forma.''}

Las fotos de moda en redes responden a un canon visual del norte global (luminosidad alta, fondos minimalistas, sonrisa estandarizada). Tu imagen propia puede repetir ese canon o atreverse a la dignidad del retrato latinoamericano: el color del Valle en tu ropa, el contexto real de tu lugar, una mirada que no necesita aprobación. La descolonización empieza por la fotografía propia.

\textbf{¿Mi imagen propia repite el canon global o se atreve a mostrar mi lugar? ¿Qué de mi contexto invisibilizo cuando elijo el fondo?}
\end{softbox}
\linead{\linewidth}\\[1mm]
\linead{\linewidth}

\begin{softbox}{milcMostaza}{milcGris}{Estoicismo · Epicteto · cuidado interior}
\textit{``No nos perturban las cosas, sino las opiniones que tenemos de ellas.''}

La trampa con la imagen propia es buscar la aprobación de los demás (likes, comentarios, comparación). Eso no depende de ti. Lo que sí depende: la decisión de transmitir dignidad y oficio, no caza de likes. Una foto que tú puedes defender es más valiosa que una con 1000 likes que no representa quién eres.

\textbf{¿Estoy tomando esta foto para mí o para perseguir la aprobación de otros? ¿Puedo defenderla en 5 años como representativa de quien era?}
\end{softbox}
\linead{\linewidth}\\[1mm]
\linead{\linewidth}

\begin{softbox}{milcTurquesa}{milcGris}{Floridi · lente de la infoesfera}
\textit{``Cada imagen publicada se vuelve dato en la infoesfera; cuídala como cuidarías un legado.''}

Tu imagen profesional, una vez publicada, alimenta sistemas de reconocimiento facial, archivos del futuro, búsquedas inversas. Es información de tu identidad biométrica. Floridi nos pediría preguntarnos: ¿qué pasa con esta imagen cuando tenga 40 años? ¿qué legado dejo en la infoesfera?

\textbf{¿Esta imagen va a envejecer bien? Si dentro de 10 años un sistema la encuentra, ¿me representará? ¿Mis decisiones de hoy honran a mi yo futuro?}
\end{softbox}
\linead{\linewidth}\\[1mm]
\linead{\linewidth}

\begin{titlebox}{milcVino}
Compromiso final
\end{titlebox}

\small
\begin{tabularx}{\textwidth}{>{\RaggedRight\arraybackslash}p{3.0cm}YYY}
\rowcolor{milcVino}\color{white}\bfseries Criterio & \color{white}\bfseries Lo logré & \color{white}\bfseries En proceso & \color{white}\bfseries Necesito apoyo\\
Comprensión del tema & Explico ideas centrales con mis palabras. & Reconozco ideas pero falta precisión. & Necesito repasar conceptos base.\\
Pensamiento computacional & Apliqué los 4 pilares al diseñar mi producto. & Apliqué algunos pilares. & Necesito releer la fase 2.\\
Aplicación y producto & Mi producto cumple los criterios y se puede revisar. & Producto funcional con ajustes pendientes. & Aún en construcción.\\
Reflexión 5 dimensiones & Respondí cada dimensión con honestidad. & Respondí algunas con detalle. & Mis respuestas son muy generales.\\
Triángulo de pensamiento & Conecté las 3 voces con mi trabajo real. & Conecté al menos una. & No relaciono las voces con mi proceso.\\
\end{tabularx}
\normalsize

\textbf{Hoy entendí que:}\\[1mm]
\linead{\linewidth}\\[1mm]
\linead{\linewidth}

\textbf{Mi compromiso para la próxima sesión es:}\\[1mm]
\linead{\linewidth}\\[1mm]
\linead{\linewidth}

\begin{softbox}{milcMostaza}{milcGris}{Cita de despedida · Marco Aurelio}
\textit{``El obstáculo en el camino se vuelve el camino. Lo que estorba al camino se vuelve camino.''}\\[1mm]
Cada error de hoy, bien observado, se convierte en pieza del camino que pisarás mañana.
\end{softbox}

\small
\begin{tabularx}{\textwidth}{>{\bfseries}p{3.6cm}Y}
\rowcolor{milcGradeDark!12}Nombre completo: & \linead{10cm}\\
Fecha: & \linead{4cm} \quad Curso/grupo: \linead{4cm}\\
Mi firma: & \linead{9cm}\\
\end{tabularx}
\normalsize

\begin{infoband}{milcGradeDark}
\textbf{Para la próxima semana.} Cambia tu foto de perfil en WhatsApp, Instagram, correo institucional y LinkedIn por tu nuevo retrato. Anota durante 3 días si te dicen algo distinto las personas con las que interactúas. Esa respuesta te dice si tu imagen funciona o necesita ajuste. El próximo lunes traes las observaciones.
\end{infoband}

\end{document}
